\subsection{Модели как инструмент проверки признаков}

В данной работе модели используются не как самостоятельный объект исследования, а как инструмент проверки информативности разных групп признаков. Основной вопрос состоит не в том, какая архитектура даёт максимальную точность, а в том, меняется ли качество прогноза при переходе от рыночных признаков к FinBERT-признакам и при их объединении.

\subsubsection{Majority baseline}

В качестве минимальной точки отсчёта используется majority baseline. Эта модель всегда предсказывает наиболее частый класс в обучающей выборке и не использует никаких признаков. Её назначение --- показать, какого качества можно достичь без рыночных или текстовых данных. Если модель на признаках не превосходит majority baseline, то её результат нельзя интерпретировать как наличие полезного прогностического сигнала.

Majority baseline особенно важен в задаче бинарного прогноза направления цены, поскольку классы могут быть несбалансированы. Например, если в определённом периоде чаще наблюдается рост, модель, всегда предсказывающая рост, может получить ненулевую accuracy без анализа признаков.

\subsubsection{Логистическая регрессия}

Основным инструментом проверки признаков является логистическая регрессия. Она выбрана как простая модель бинарной классификации, позволяющая оценить базовую информативность признаков без усложнения архитектуры. Логистическая регрессия обучается в трёх вариантах:

\begin{enumerate}
    \item \texttt{prices-only} --- только рыночные признаки;
    \item \texttt{FinBERT-only} --- только embedding-признаки финансовых текстов;
    \item \texttt{prices + FinBERT} --- объединение рыночных и текстовых признаков.
\end{enumerate}

Такое сравнение позволяет разделить три эффекта. Модель \texttt{prices-only} показывает вклад исторической ценовой динамики. Модель \texttt{FinBERT-only} проверяет, есть ли самостоятельный сигнал в текстах. Модель \texttt{prices + FinBERT} проверяет, дают ли текстовые признаки добавочную силу поверх рыночных.

Для логистической регрессии используется балансировка классов: \texttt{class\_weight = balanced}. Это снижает риск того, что модель будет ориентироваться только на более частый класс. Однако итоговая интерпретация строится не только по accuracy, но и по balanced accuracy, F1-score, ROC-AUC и confusion matrix.

\subsubsection{LSTM}

Дополнительно используется LSTM-модель \cite{hochreiter1997lstm}. Её роль в работе --- проверить, меняется ли эффект признаков при использовании модели, способной учитывать последовательную структуру данных. В отличие от логистической регрессии, LSTM получает на вход окно из нескольких предыдущих дней и может учитывать порядок наблюдений внутри этого окна.

LSTM применяется к объединённым признакам \texttt{prices + FinBERT}. Такой вариант выбран потому, что именно объединённая группа является основной для проверки гипотезы о добавочной силе текстовой информации. Если последовательная модель извлекает из объединённых признаков больше сигнала, чем логистическая регрессия, это может указывать на важность временной структуры. Если же LSTM не улучшает качество, то усложнение модели не подтверждает дополнительную полезность признаков в текущем протоколе.

Архитектурно модель принимает последовательность дневных признаков и возвращает бинарный прогноз направления движения цены. Для обучения используется \texttt{BCEWithLogitsLoss}, а качество контролируется на validation subset. Early stopping применяется для ограничения переобучения: модель не выбирается по минимальной ошибке на train subset, а останавливается при отсутствии улучшения на validation.

Таким образом, модели в данной работе выполняют вспомогательную функцию. Они нужны для эмпирической проверки того, какие группы признаков связаны с последующим движением цены и усиливает ли текстовая информация прогноз относительно рыночной динамики.
